% from ega4-20-fr.tex
\oldpage[IV-4]{242}
  dans son adhérence, ce qui prouve déjà que l'espace sous-jacent à $W$ est ouvert dans $T$ ; la conclusion résulte alors de \sref{IV.11.10.1}[(c)].

  Ce lemme étant établi, on en conclut que pour $\mu\geq\lambda$ assez grand, $Z_\mu$ est induit sur un ouvert de $X_\mu$ en vertu de \sref{IV.8.6.3}, car $Z_\lambda$, en tant que sous-préschéma d'un préschéma noethérien, est de présentation finie sur $X_\lambda$ \sref{IV.1.6.2}, et il en est donc de même des $Z_\mu$ pour $\mu\geq\lambda$ sur $X_\mu$ et de $Z$ sur $X$.
  On peut maintenant appliquer \sref{IV.20.3.8.2} qui montre que pour $\mu\geq\lambda$ assez grand, $Z_\mu$ est schématiquement dense dans $X_\mu$, d'où la conclusion.

  \item Si, dans l'hypothèse~(i) de \sref{IV.20.3.8}, on supprime la condition que $X$ est noethérien, on voit que le raisonnement de \sref{IV.20.3.8.6} montre encore que l'image de \sref{IV.20.3.8.1} est formée des pseudo-$S$-morphismes ayant un représentant qui est un $S$-morphisme $U\to Y$, où $U$ est schématiquement dense dans $X$ et \emph{quasi-compact et quasi-séparé}.
\end{enumerate}
\end{remarks}
\end{proof}

\begin{corollary}[20.3.9]
\label{IV.20.3.9}
Supposons vérifiées l'une ou l'autre hypothèse~a),~b) de \sref{IV.20.3.6} sur $X$, et que $Y$ soit localement de présentation finie sur $S$.
Alors, pour qu'un pseudo-$S$-morphisme $\omega$ de $X$ dans $Y$ soit défini au point $x$ \sref{IV.20.2.3}, il faut et il suffit que sa restriction à $\Spec(\sh{O}_{X,x})$ soit partout définie (autrement dit, soit un $S$-morphisme de $\Spec(\sh{O}_{X,x})$ dans $Y$).
\end{corollary}

Le résultat suivant, qui nous servira dans la démonstration de \sref{IV.20.3.11}, utilise la théorie de la descente fidèlement plate du chap.~VI.
Le lecteur pourra vérifier que les résultats de \sref{IV.20.3} ne seront pas utilisés dans cette théorie.

\begin{lemma}[20.3.10]
\label{IV.20.3.10}
Soient $f:X'\to X$ un $S$-morphisme fidèlement plat et quasi-compact, $X''=X'\times_X X'$, $p_1$ et $p_2$ les projections canoniques de $X''$ dans $X'$, $Y$ un préschéma séparé sur $S$.
Soient $U$ un ouvert de $X$, $U'=f^{-1}(U)$, $U''=p_1^{-1}(U')=p_2^{-1}(U')$, et supposons que $U''$ soit schématiquement dense dans $X''$.
Soit $g:U\to Y$ un $S$-morphisme ; alors, si $g\circ(f|U')$ se prolonge en un $S$-morphisme $X'\to Y$, $g$ se prolonge en un $S$-morphisme $X\to Y$.
\end{lemma}

On notera que les hypothèses entraînent que $U$ (resp. $U'$) est schématiquement dense dans $X$ (resp. $X'$) \sref{IV.11.10.5}[(i)] ; on peut donc encore dire que si $\omega$ désigne le pseudo-$S$-morphisme classe de $g$, l'énoncé de \sref{IV.20.3.10} signifie que \emph{si $\omega\circ f$ est partout défini, il en est de même de $\omega$}.

\begin{proof}
Pour prouver \sref{IV.20.3.10}, désignons par $g'$ un $S$-morphisme $X'\to Y$ qui prolonge $g\circ(f|U')$, et posons $g''_i=g'\circ p_i:X''\to Y$ ($i=1,2$).
Si l'on pose $f''=f\circ p_1=f\circ p_2:X''\to X$, il est clair que $g''_1$ et $g''_2$ coïncident dans $U''$ avec $g\circ(f''|U'')$.
Mais comme $Y$ est séparé sur $S$ et $U''$ schématiquement dense dans $X''$, on a $g''_1=g''_2$ \sref{IV.11.10.1}[(d)].
Comme $f$ est fidèlement plat et quasi-compact, il résulte de la théorie de la descente (chap.~VI) qu'il existe un unique $S$-morphisme $h:X\to Y$ tel que $h\circ f=g'$ ; comme la restriction $U'\to U$ de $f$ est un morphisme fidèlement plat et quasi-compact et que $U''=U'\times_U U'$, le résultat d'unicité précédent, appliqué à $U$ au lieu de $X$, montre que $h|U=g$, ce qui prouve le lemme.
\end{proof}

\begin{proposition}[20.3.11]
\label{IV.20.3.11}
Soient $Y$ un $S$-préschéma séparé sur $S$, $\omega$ un pseudo-$S$-morphisme de $X$ dans $Y$, $f:X'\to X$ un $S$-morphisme.
Supposons que $f$ soit plat, et que l'une des conditions suivantes soit réalisée :
\begin{enumerate}[label=\emph{(\roman*)}]
